\scanpage{15}
\[
\unsure{n}\longmapsto\mathrm{rg}_K L_n\otimes_k K
\]
est polynômiale \deleted{pour $n$ grand}. [le degré du polynôme \inserted{(supposé non nul, i.e. $L_\bullet\ne0$)} est alors le plus grand entier $\unsure{+n|n|N}$ tel que $L_n^!\ne0$, et les coefficients du polynôme exprimé comme combinaison linéaire des $P_i(n)=\binom n i$ $(i\geq0)$ sont les $\mathrm{rg}_K L_i^!\otimes_k K$ ; les coefficients dudit polynôme (par r. aux $P_i$, ou pour l'écriture habituelle \deleted{et} ou comme comb. lin. des \unsure{$n\mapsto n^i$}) sont des fonctions constructibles en le pt variable $\unsure{\mathfrak p|p}\in\operatorname{Spec} k$].

\uline{NB} Si dans e) on suppose $k$ noeth mais non régulier, la condition \inserted{\unsure{énoncée}} de $\partial$-perfection est tjrs néc, mais pas suffisante -- il faut une condition néc. et suff. plus subtile, savoir l'existence d'un $L'_\bullet\unsure{\longrightarrow|\mathrel{-}}L_\bullet$, avec $L'_\bullet$ \uline{str.} $\partial$-parfait (cf c)) et que \deleted{soit un \unsure{isom.}} \inserted{\unsure{qu-i|quis}} i.e. induise un iso sur les $\pi_i$. (De plus, \unsure{$\exists\exists$|$\exists$} $L'\to L$ épi)

% slanted
\margin{Stabilité des notions \unsure{préc.} pour sommes, \unsure{extensions} ; \unsure{et pour} produits \unsure{tensoriels} $\Lambda^i,\Gamma^i,\operatorname{Sym}^i$ ! dans le cas \uuline{str.} $\partial$ ps. coh. et \uuline{str.} $\partial$ parfaits \deleted{\unsure{c'est évident}}}

On veut un théorème : \uline{\unsure{Pour|Dévis.}} \uline{\unsure{tout|précis}} des $L_\bullet$ \unsure{str. $\partial$-ps. coh.} \deleted{\unsure{on}} -- cela n'est possible bien sûr que \deleted{plus subtilt de \unsure{ps.} $\partial$-ps. coh.}
% archive sheet number: 14
